%% Requires compilation with XeLaTeX or LuaLaTeX
\documentclass[10pt,xcolor={table,dvipsnames},t]{beamer}
\usetheme{UVic}

\usepackage{booktabs}

\title[PhD Progress Report]{PhD Progress Report}
\subtitle{Learning Automata Framework for Adaptive Kelp Forest Monitoring in British Columbia Coastal Waters}
\author{Ekaba Bisong}
\institute{Department of Computer Science\\University of Victoria}
\date{December 2025}

\begin{document}

% Slide 1: Title
\begin{frame}
  \titlepage
\end{frame}

% Slide 2: Research Overview
\begin{frame}{Research Overview}

\textbf{Problem:} Monitoring kelp forests via satellite imagery is challenging due to:
\begin{itemize}
  \item Limited labeled training data
  \item Heterogeneous environmental conditions along BC's coastline
  \item Sparse and delayed ground-truth validation
\end{itemize}

\vspace{0.5em}
\textbf{Solution:} An adaptive Learning Automata (LA) framework that:
\begin{itemize}
  \item Learns to select optimal detection heuristics with minimal supervision
  \item Adapts to varying turbidity, SST, and seasonal patterns
  \item Handles asynchronous label arrival through buffering mechanisms
\end{itemize}

\end{frame}

% Slide 3: Main Objectives
\begin{frame}{Main Objectives}

\begin{enumerate}
  \item \textbf{Adaptive LA Framework:} Develop a Learning Automata framework for satellite-based kelp monitoring with limited labeled data

  \item \textbf{Context-Aware Selection:} Design heuristic selection mechanisms that adapt to heterogeneous environmental conditions

  \item \textbf{Sparse Label Handling:} Create buffering and delayed reward propagation for asynchronous ground-truth data

  \item \textbf{TEK Integration:} Establish protocols for integrating Traditional Ecological Knowledge while respecting indigenous data sovereignty

  \item \textbf{Experimental Design:} Formulate complete evaluation using BC Bull Kelp dataset with ablation studies
\end{enumerate}

\end{frame}

% Slide 4: Completed Work
\begin{frame}{Completed Work}

\textbf{Thesis Proposal Complete:}
\begin{itemize}
  \item Comprehensive literature review (kelp ecology, remote sensing, LA theory)
  \item Core methodology with three integrated components:
  \begin{itemize}
    \item Component 1: Global Policy Learning (L$_{\text{R-I}}$ updates)
    \item Component 2: Context-Aware Local Adaptation (memory banks, k-NN)
    \item Component 3: Sparse Label Handling (temporal decay)
  \end{itemize}
  \item Conceptual models for all four research questions (RQ1--RQ4)
  \item Experimental design with stratified evaluation protocols
  \item TEK integration protocol with phased community engagement
\end{itemize}

\end{frame}

% Slide 5: Research Questions
\begin{frame}{Research Questions}

\begin{itemize}
  \item \textbf{RQ1:} How effectively can an LA framework adapt to new regions with limited labeled examples?

  \vspace{0.5em}
  \item \textbf{RQ2:} Can context-aware heuristic selection improve detection across heterogeneous coastal environments?

  \vspace{0.5em}
  \item \textbf{RQ3:} How can delayed and sparse feedback be incorporated into the learning process?

  \vspace{0.5em}
  \item \textbf{RQ4:} What protocols enable respectful integration of Traditional Ecological Knowledge?
\end{itemize}

\end{frame}

% Slide 6: Timeline 2026
\begin{frame}{Timeline: 2026}

\begin{itemize}
  \item \textbf{January 2026} --- Proposal Defense: Committee approves PhD proposal

  \vspace{0.3em}
  \item \textbf{May 2026} --- M1: Dataset and baselines ready
  \begin{itemize}
    \item Sentinel-2 imagery preprocessed
    \item Ground-truth labels compiled
    \item Tier 1 spectral indices implemented
  \end{itemize}

  \vspace{0.3em}
  \item \textbf{October 2026} --- M2: Core framework functional
  \begin{itemize}
    \item All three LA components integrated
    \item Tier 2 and Tier 3 heuristics developed
  \end{itemize}

  \vspace{0.3em}
  \item \textbf{November 2026} --- \textbf{Paper 1 (RQ1):} ``Learning Automata for Adaptive Kelp Detection with Limited Labels''
\end{itemize}

\end{frame}

% Slide 7: Timeline 2027
\begin{frame}{Timeline: 2027}

\begin{itemize}
  \item \textbf{January 2027} --- M3: Temporal adaptation experiments complete

  \item \textbf{February 2027} --- \textbf{Paper 2 (RQ2):} ``Context-Aware Heuristic Selection for Heterogeneous Coastal Environments''

  \item \textbf{March 2027} --- M4 \& M5: Spatial and sparse label experiments complete

  \item \textbf{April 2027} --- \textbf{Paper 3 (RQ3):} ``Learning from Delayed and Sparse Feedback in Environmental Monitoring''

  \item \textbf{May 2027} --- M6: Extensions implemented (pursuit algorithms, TEK Phase 1)

  \item \textbf{June 2027} --- \textbf{Paper 4 (RQ4):} ``Integrating Traditional Ecological Knowledge into Automated Kelp Forest Monitoring''

  \item \textbf{August 2027} --- M7: Thesis draft complete

  \item \textbf{October 2027} --- M8: Thesis defense
\end{itemize}

\end{frame}

% Slide 8: Publication Plan
\begin{frame}{Publication Plan}

\textbf{Four Papers Aligned with Research Questions:}

\vspace{0.5em}
\begin{itemize}
  \item \textbf{Paper 1 (RQ1)} --- Core LA framework, adaptation with minimal supervision
  \begin{itemize}
    \item Target: November 2026
  \end{itemize}

  \item \textbf{Paper 2 (RQ2)} --- Context-aware heuristic selection across varying conditions
  \begin{itemize}
    \item Target: February 2027
  \end{itemize}

  \item \textbf{Paper 3 (RQ3)} --- Buffering mechanisms, temporal decay, asynchronous labels
  \begin{itemize}
    \item Target: April 2027
  \end{itemize}

  \item \textbf{Paper 4 (RQ4)} --- TEK integration and indigenous data sovereignty
  \begin{itemize}
    \item Target: June 2027
  \end{itemize}
\end{itemize}

\end{frame}

% Slide 9: Next Steps
\begin{frame}{Immediate Next Steps}

\textbf{Next 2--3 Months:}
\begin{itemize}
  \item Finalize PhD proposal for committee review
  \item Proposal defense (January 2026)
  \item Begin Phase 1: Foundation and Data Preparation
  \begin{itemize}
    \item Acquire and preprocess Sentinel-2 imagery for BC coastal waters
    \item Compile ground-truth labels from available sources
    \item Implement baseline spectral indices (NDVI, FAI, NDWI, Kelp Index)
    \item Establish evaluation infrastructure
  \end{itemize}
\end{itemize}

\end{frame}

% Slide 10: Summary
\begin{frame}{Summary}

\begin{itemize}
  \item \textbf{Proposal Status:} Complete and ready for committee review

  \vspace{0.5em}
  \item \textbf{Methodology:} Three-component LA framework designed with mathematical formalization

  \vspace{0.5em}
  % \item \textbf{Timeline:} On track for October 2027 completion

  \vspace{0.5em}
  \item \textbf{Publications:} Four papers planned, one per research question

  \vspace{0.5em}
  \item \textbf{Next Milestone:} Proposal defense (January 2026)
\end{itemize}

\vspace{1em}
\centering
\textbf{Questions?}

\end{frame}

\end{document}
